\chapter{Improper integrals, convergence, and engineered endpoints}

\begin{orient}
\textbf{What this chapter is for.} Every integral in this chapter asks two questions, and they must be answered in that order. The first is whether the integral exists at all. The second is what it equals. Answering the second without the first is the commonest way to write a confident wrong answer in all of integration, because the machinery of antidifferentiation is perfectly happy to produce a finite number for a divergent object: the fundamental theorem applied across a pole gives $\int_{-1}^{1}x^{-2}\dd x=-2$, a value that is not merely inaccurate but meaningless, and nothing in the algebra warns you. So the chapter begins with a taxonomy of improperness (infinite interval, unbounded integrand, both at once, and the interior singularity that hides in plain sight), then builds the convergence machinery: comparison, limit comparison, the $p$-test at zero and at infinity, and the separate regime of oscillatory integrands that converge without converging absolutely. It then treats the two respectable ways to attach a number to a divergent object, the Cauchy principal value and the cancellation of a matched pair of logarithmic divergences, and it insists that both be labelled as what they are. The chapter closes on a different skill entirely: reading engineered endpoints. When a problem hands you a limit of integration written as $\sqrt{\eu^{2}-1}$ or $\log_{8}\eu$, that number is a message about the intended substitution, and decoding it is faster than any amount of pattern-matching on the integrand.
\end{orient}

\section{Two questions, in order}

Write $\int_a^bf(x)\dd x$ and the notation quietly asserts three things: that $[a,b]$ is a bounded interval, that $f$ is bounded on it, and that the Riemann or Lebesgue construction therefore applies directly. An integral is \emph{improper} when one of those assertions fails. It is then not an integral at all until you have defined it as a limit, and the definition is a genuine part of the answer rather than a formality to be waved at.

There are exactly three ways an assertion can fail, and a fourth case that is a combination.

\begin{center}
\footnotesize
\setlength{\tabcolsep}{5pt}
\renewcommand{\arraystretch}{1.32}
\begin{tabular}{@{}llll@{}}
\toprule
\mh{Type} & \mh{Symptom} & \mh{Definition} & \mh{Model case} \\
\midrule
I & Infinite bound & $\displaystyle\lim_{R\to\infty}\int_a^{R}f$ & $\displaystyle\int_1^{\infty}x^{-p}\dd x$ \\
II & Endpoint singularity & $\displaystyle\lim_{\varepsilon\to0^{+}}\int_{a+\varepsilon}^{b}f$ & $\displaystyle\int_0^{1}x^{-p}\dd x$ \\
III & Interior singularity at $c$ & $\displaystyle\int_a^{c}+\int_c^{b}$, both required & $\displaystyle\int_{-1}^{1}x^{-2}\dd x$ \\
Mixed & Both ends offend & Split at any interior point & $\displaystyle\int_0^{\infty}\frac{x^{s-1}}{1+x^{n}}\dd x$ \\
\bottomrule
\end{tabular}
\end{center}

The table also settles a question of notation that causes real confusion. An improper integral of Type I and an improper integral of Type II are the same object viewed through a substitution: $x=1/u$ carries $\int_1^{\infty}x^{-p}\dd x$ to $\int_0^{1}u^{p-2}\dd u$, and the two $p$-tests are one $p$-test seen from two sides. This is more than an observation. It is a working technique, because a singularity at $0$ is often easier to analyse after being pushed out to $\infty$ by $x=\eu^{-t}$ or $x=1/u$, and an infinite range is often easier to integrate numerically after being pulled in to a finite one by the same map.

Type III is the dangerous one, because it is the only type with no visible marker. An infinite bound announces itself. An endpoint singularity is usually obvious from the presence of $\ln x$ at $0$ or $\sqrt{1-x^{2}}$ downstairs at $1$. An interior pole is invisible unless you look for it: $\int_0^{3}\dfrac{\dd x}{x-1}$ looks like an ordinary rational integral, and its integrand is finite at both endpoints. The habit that prevents this failure is mechanical. Before writing an antiderivative, solve ``denominator $=0$'' and ``argument of $\ln$ or of an even root $=0$'' on the closed interval, and check whether any solution lies inside.

\begin{keynote}[The two questions]
\textbf{Does it converge?} is a question about the integrand's behaviour near the offending points only, and it is answered by comparison with a power. \textbf{What is it?} is a question about the whole interval, and it is answered by the techniques of every other chapter. The first question is cheap: it takes one line per bad point. The second is expensive. Doing them out of order means paying the expensive one before knowing whether it has an answer, and worse, it means that a divergent integral will present itself to you as a finite number with no indication that anything went wrong.
\end{keynote}

The mixed case deserves a word because it is the standard shape of the integrals that matter. An integral over $(0,\infty)$ of a rational or algebraic integrand is almost always improper at both ends, and the two ends impose \emph{independent} conditions. The archetype is $\int_0^{\infty}\dfrac{x^{s-1}}{1+x^{n}}\dd x$: near $0$ the integrand behaves like $x^{s-1}$, which needs $s>0$; near $\infty$ it behaves like $x^{s-1-n}$, which needs $s-1-n<-1$, that is $s<n$. The verdict is the conjunction, $0<s<n$, and a candidate closed form asserted outside that strip is not a value but a formal expression. This is why the convergence step is not pedantry: it is what tells you the domain on which your answer is an answer.

\section{Endpoints as limits}

The first technique is a discipline rather than a trick, and it is the one on which everything else in the chapter rests. Its content is that the offending endpoint must be replaced by a variable and the limit taken at the end, after the antiderivative has been evaluated symbolically.

\tech{Endpoint limit discipline}{core}
\cue An infinite bound; an integrand undefined at an endpoint ($\ln x$ at $0$, $\dfrac{1}{\sqrt{1-x^{2}}}$ at $1$, $x^{-1/2}$ at $0$); or a zero of a denominator strictly inside the interval.
\method Replace every offending endpoint by a limit \emph{before} manipulating:
\[\int_a^bf(x)\dd x=\lim_{\varepsilon\to0^{+}}\int_{a+\varepsilon}^{b}f(x)\dd x,\qquad \int_a^{\infty}f(x)\dd x=\lim_{R\to\infty}\int_a^{R}f(x)\dd x.\]
For a singularity at an interior point $c$, split into $\int_a^{c}+\int_c^{b}$ and require \emph{both} pieces to converge independently, with independent cut-offs. Before splitting, decide whether the singularity is genuine: if $f$ extends continuously across $c$ (as $\dfrac{\sin x}{x}$ does at $0$, or $\dfrac{\ln x}{1-x}$ at $1$), the integral is proper and no split is needed.

\begin{worked}
\[\int_0^{1}\ln x\dd x=\lim_{\varepsilon\to0^{+}}\bigl[x\ln x-x\bigr]_{\varepsilon}^{1}=\lim_{\varepsilon\to0^{+}}\bigl(-1-\varepsilon\ln\varepsilon+\varepsilon\bigr)=-1,\]
because $\varepsilon\ln\varepsilon\to0$. The limit is not decoration: the antiderivative $x\ln x-x$ is undefined at $x=0$, and the value $-1$ exists only as a limit of the truncated integrals.

A mixed case, both ends offending:
\[\int_0^{\infty}\frac{\dd x}{\sqrt{x}\,(1+x)}.\]
Near $0$ the integrand is $\sim x^{-1/2}$, integrable; near $\infty$ it is $\sim x^{-3/2}$, integrable. So the object exists. Now evaluate: $x=t^{2}$ gives $\int_0^{\infty}\dfrac{2\dd t}{1+t^{2}}=\pi\approx3.141593$. Two lines of convergence, one line of evaluation.
\end{worked}

\begin{worked}[A removable singularity, correctly not split]
\[\int_0^{1}\frac{\ln x}{1-x}\dd x.\]
The denominator vanishes at $x=1$, so the reflex is to split. But $\ln x=\ln(1-(1-x))\sim-(1-x)$ near $1$, so the integrand tends to $-1$ there and the point is removable. At the other end $\ln x$ has an integrable singularity. Hence the integral is improper only at $0$ and converges. Expanding $\dfrac{1}{1-x}=\sum_{n\geq0}x^{n}$ and integrating termwise gives $\sum_{n\geq0}\dfrac{-1}{(n+1)^{2}}=-\dfrac{\pi^{2}}{6}\approx-1.644934$.
\end{worked}

\begin{worked}[An interior singularity handled correctly]
\[\int_0^{2}\frac{\dd x}{\sqrt{\abs{x-1}}}.\]
The integrand blows up at the interior point $x=1$, so the object is Type III and the fundamental theorem may not be applied across it. Split and take two independent limits:
\[\lim_{\varepsilon\to0^{+}}\int_0^{1-\varepsilon}\frac{\dd x}{\sqrt{1-x}}=\lim_{\varepsilon\to0^{+}}\bigl[-2\sqrt{1-x}\bigr]_0^{1-\varepsilon}=2,\qquad \lim_{\delta\to0^{+}}\int_{1+\delta}^{2}\frac{\dd x}{\sqrt{x-1}}=2,\]
and both exist, so the integral converges to $4$. The exponent is what saves it: $p=\tfrac12<1$ on each side. Replace the square root by a first power and both halves become $+\infty$, so the integral diverges even though the picture looks almost identical.
\end{worked}

\begin{trap}
Applying the fundamental theorem across an interior pole. The antiderivative of $x^{-2}$ is $-x^{-1}$, and $\bigl[-x^{-1}\bigr]_{-1}^{1}=-1-1=-2$, a negative number for a strictly positive integrand: the sign alone should convict it. The integral diverges, since $\int_0^{1}x^{-2}\dd x=+\infty$. The second form of the same error is splitting at $c$ but tying the two cut-offs together, writing $\lim_{\varepsilon\to0}\bigl[\int_a^{c-\varepsilon}+\int_{c+\varepsilon}^{b}\bigr]$ and calling the result the integral. That limit is the principal value, which is a different and weaker object; the integral requires the two limits to exist \emph{separately}.
\end{trap}

\begin{seealsob}Comparison and limit comparison; Cauchy principal value; splitting at a symmetry or sign-change point; parity on a symmetric interval.\end{seealsob}

The discipline tells you to take a limit. It does not tell you whether the limit exists, and for anything but the simplest integrands you cannot find out by computing the antiderivative, because there may not be one. What is needed is a test that inspects only the local behaviour at the bad point, and that is exactly what comparison provides.

\section{Deciding convergence}

Convergence of an improper integral is decided entirely by the integrand's rate of blow-up or decay near each offending point, and rates are measured against powers. Two facts carry almost all of the load, and they point in opposite directions, which is why they are so often confused.

\begin{keynote}[The $p$-test, both ends]
\[\int_1^{\infty}\frac{\dd x}{x^{p}}\ \text{converges}\iff p>1,\qquad \int_0^{1}\frac{\dd x}{x^{p}}\ \text{converges}\iff p<1.\]
At infinity you need \emph{fast} decay; at the origin you need \emph{mild} blow-up. The exponent $p=1$ fails at both ends, and $\int_0^{\infty}\dfrac{\dd x}{x}$ diverges twice over. Logarithms are weaker than any power and therefore do not move the verdict except exactly at $p=1$, where they decide it: $\int_2^{\infty}\dfrac{\dd x}{x\ln^{q}x}$ converges iff $q>1$, while $\int_1^{\infty}\dfrac{\ln^{k}x}{x^{p}}\dd x$ converges iff $p>1$ for every fixed $k$. Near the origin the corresponding statement is that $x^{-p}\ln^{k}x$ is integrable for every $k$ when $p<1$, and for no $k$ when $p\geq1$.
\end{keynote}

The reason the $p$-test suffices is that it is applied through comparison, and comparison only needs an inequality near the bad point. Everything away from the singularity is a proper integral of a bounded function and contributes a finite amount that cannot change the verdict. It follows that convergence is a purely local property: you may modify $f$ arbitrarily on any interval bounded away from the bad points without changing the answer to the first question, and that is why the analysis reduces to matching a power at each of at most two places.

A small reference library makes the matching instantaneous. Every entry below is worth knowing as a fact rather than rederiving, and every one of them is the comparison function for a large class of problems.

\begin{center}
\footnotesize
\setlength{\tabcolsep}{5pt}
\renewcommand{\arraystretch}{1.32}
\begin{tabular}{@{}lll@{}}
\toprule
\mh{Integral} & \mh{Verdict} & \mh{Value when convergent} \\
\midrule
$\displaystyle\int_0^{1}x^{-p}\dd x$ & $p<1$ & $\dfrac{1}{1-p}$ \\
$\displaystyle\int_1^{\infty}x^{-p}\dd x$ & $p>1$ & $\dfrac{1}{p-1}$ \\
$\displaystyle\int_0^{1}\ln^{k}x\dd x$ & always & $(-1)^{k}k!$ \\
$\displaystyle\int_0^{1/2}\frac{\dd x}{x\ln^{2}x}$ & converges & $\dfrac{1}{\ln2}\approx1.442695$ \\
$\displaystyle\int_2^{\infty}\frac{\dd x}{x\ln^{q}x}$ & $q>1$ & $\dfrac{\ln^{1-q}2}{q-1}$ \\
$\displaystyle\int_0^{\infty}x^{s-1}\eu^{-x}\dd x$ & $s>0$ & $\Gfun(s)$ \\
$\displaystyle\int_0^{1}\frac{\dd x}{\sqrt{1-x^{2}}}$ & converges & $\dfrac{\pi}{2}$ \\
$\displaystyle\int_0^{\infty}\frac{\sin x}{x}\dd x$ & conditional & $\dfrac{\pi}{2}$ \\
\bottomrule
\end{tabular}
\end{center}

The fourth row repays a second look, because it is the one that surprises. Near the origin $\dfrac{1}{x\ln^{2}x}$ blows up like $x^{-1}$ up to a logarithmic factor, and $x^{-1}$ diverges; yet the squared logarithm downstairs is enough to rescue it, since the substitution $u=\ln x$ turns the integral into $\int_{-\infty}^{-\ln2}u^{-2}\dd u$, which converges at $-\infty$. The general principle: at the exact threshold $p=1$ the logarithms decide, and they decide by the same $q>1$ rule one level down.

\tech[Comparison and limit comparison for integrals]{Comparison and limit comparison}{core}
\cue You are about to assert a value for an integral over an unbounded domain, or with an unbounded integrand, and you have not yet checked that the object exists.
\method Isolate each bad point. Near it, find the dominant behaviour of $\abs{f}$ and compare with a power. Direct comparison: if $0\leq f\leq g$ near the bad point and $\int g$ converges, so does $\int f$; if $f\geq g\geq0$ and $\int g$ diverges, so does $\int f$. Limit comparison is usually faster and needs no inequality:
\[\lim_{x\to c}\frac{f(x)}{g(x)}=L\in(0,\infty)\ \Longrightarrow\ \int f\ \text{and}\ \int g\ \text{share their verdict}.\]
Choose $g$ to be the power that matches the leading behaviour: strip off every factor that tends to a nonzero constant. Then test each bad point separately and take the conjunction.

\begin{worked}
\[\int_0^{\infty}\frac{x^{s-1}}{1+x^{n}}\dd x,\qquad n>0.\]
Two bad points. At $0$: $\dfrac{x^{s-1}}{1+x^{n}}\big/x^{s-1}\to1$, so the piece $\int_0^{1}$ behaves like $\int_0^{1}x^{s-1}\dd x$, convergent iff $1-s<1$, that is $s>0$. At $\infty$: $\dfrac{x^{s-1}}{1+x^{n}}\big/x^{s-1-n}\to1$, so $\int_1^{\infty}$ behaves like $\int_1^{\infty}x^{s-1-n}\dd x$, convergent iff $n-s+1>1$, that is $s<n$. The integral converges precisely for $0<s<n$, and on that strip its value is $\dfrac{\pi}{n\sin(\pi s/n)}$. At $s=\tfrac32$, $n=4$ this reads $\dfrac{\pi}{4\sin(3\pi/8)}\approx0.850109$, matching quadrature to ten digits.
\end{worked}

\begin{worked}[Choosing the comparison power by inspection]
$\displaystyle\int_1^{\infty}\frac{\dd x}{\sqrt{x^{3}+1}}$: the radicand is dominated by $x^{3}$, so the integrand is $\sim x^{-3/2}$ and the limit-comparison ratio tends to $1$. Since $\tfrac32>1$, it converges.

$\displaystyle\int_0^{1}\frac{\dd x}{\sqrt{x}-x}$: at $0$ the dominant term downstairs is $\sqrt x$, giving $\sim x^{-1/2}$ with $p=\tfrac12<1$, convergent. But there is a second bad point: the denominator also vanishes at $x=1$, where $\sqrt x-x=\sqrt x(1-\sqrt x)\sim\tfrac12(1-x)$, giving $p=1$ and divergence. The integral diverges, and it diverges at the end nobody looked at.

$\displaystyle\int_1^{\infty}x^{100}\eu^{-\sqrt x}\dd x$: no power test applies directly, and none is needed. Exponential decay beats every power, so compare with $x^{-2}$: the ratio $x^{102}\eu^{-\sqrt x}\to0$, which gives $f\leq x^{-2}$ eventually and convergence. The rule to internalise is the hierarchy $\ln^{k}x\ll x^{a}\ll\eu^{bx}$ for all $k$ and all $a,b>0$; against an exponential, every polynomial and logarithmic decoration is invisible to the convergence question.
\end{worked}

\begin{trap}
Testing one endpoint and forgetting the other. A power that is comfortable at $\infty$ is usually the one that kills you at $0$, precisely because the two conditions run in opposite directions, and the second worked example above shows the failure in its purest form: the visible singularity converges and the invisible one does not. The second habitual error is comparing $\abs{f}$ when $f$ changes sign and then concluding \emph{divergence} from the divergence of $\int\abs f$. Comparison proves convergence of $\int f$ from convergence of $\int\abs f$, never the reverse.
\end{trap}

\begin{seealsob}Endpoint limit discipline; oscillatory convergence and Dirichlet's test; Euler reflection and the Mellin transform.\end{seealsob}

Comparison is a statement about $\abs f$, so it establishes \emph{absolute} convergence, and absolute convergence is what every interchange theorem in analysis wants: dominated convergence, Tonelli, termwise integration of a series, differentiation under the integral sign. When $\int\abs f$ diverges but $\int f$ converges anyway, the integral still has a value, but it has lost the right to be manipulated freely, and that distinction is the subject of the next technique.

\section{Absolute versus conditional}

An integral converges absolutely when $\int\abs f$ converges. It converges conditionally when $\int f$ converges and $\int\abs f$ does not. The prototype is $\int_0^{\infty}\dfrac{\sin x}{x}\dd x$, which equals $\dfrac{\pi}{2}$ while $\int_0^{\infty}\Bigl|\dfrac{\sin x}{x}\Bigr|\dd x=\infty$: on the $n$-th half-period the absolute integrand contributes at least $\dfrac{2}{\pi(n+1)}$, and the harmonic series diverges. The convergence of the signed integral comes from cancellation between successive humps, and cancellation is fragile. Rearranging it, bounding it, or passing a limit through it can destroy it.

\tech[Oscillatory convergence and Dirichlets test]{Oscillatory convergence and Dirichlet's test}{hard}
\cue $\dfrac{\sin x}{x}$, $\dfrac{\cos x}{\sqrt x}$, $\sin(x^{2})$, $\dfrac{\sin x}{\ln x}$: an integrand that decays too slowly for absolute convergence but oscillates with a bounded primitive.
\method Dirichlet's test. If $g$ is monotone on $[a,\infty)$ with $g(x)\to0$, and the partial integrals $\int_a^{X}f$ are bounded uniformly in $X$, then $\int_a^{\infty}f g$ converges. In practice there are two working routes. Integrate by parts once to move the decay into the denominator, converting a conditionally convergent integral into an absolutely convergent one; or split at the zeros of the oscillating factor and treat the result as an alternating series with decreasing terms.

\begin{worked}
$\displaystyle\int_1^{\infty}\frac{\sin x}{x}\dd x$. Take $f=\sin x$, whose partial integrals $\bigl[-\cos x\bigr]_1^{X}$ are bounded by $2$, and $g=\tfrac1x$, monotone to $0$. Dirichlet gives convergence. Integrating by parts exhibits the mechanism:
\[\int_1^{R}\frac{\sin x}{x}\dd x=\Bigl[\frac{-\cos x}{x}\Bigr]_1^{R}-\int_1^{R}\frac{\cos x}{x^{2}}\dd x,\]
and the remaining integrand is $O(x^{-2})$, absolutely convergent. Together with $\int_0^{1}\tfrac{\sin x}{x}\dd x$, which is proper because the singularity is removable, the total is $\dfrac{\pi}{2}\approx1.570796$.
\end{worked}

\begin{worked}[Oscillation manufactured by a substitution]
$\displaystyle\int_0^{\infty}\sin(x^{2})\dd x$ has an integrand that does not decay at all. Substitute $u=x^{2}$:
\[\int_0^{\infty}\sin(x^{2})\dd x=\int_0^{\infty}\frac{\sin u}{2\sqrt u}\dd u,\]
which is exactly the Dirichlet shape with $g(u)=\tfrac{1}{2\sqrt u}$ monotone to $0$. It converges conditionally, to $\tfrac12\sqrt{\tfrac{\pi}{2}}\approx0.626657$. The decay was hidden in the Jacobian, not in the integrand: the reason $\sin(x^{2})$ converges is that its oscillations become infinitely fast, so each hump has vanishing width.
\end{worked}

\begin{trap}
Applying dominated convergence, Tonelli, or an absolute comparison to a conditionally convergent integral. Every one of those theorems has an integrable dominating function in its hypotheses, and no such function exists here. Concretely: swapping $\int_0^{\infty}$ with $\sum$ in $\int_0^{\infty}\tfrac{\sin x}{x}\dd x$ after expanding $\tfrac1x$ in some series, or interchanging the order in a double integral to derive $\tfrac{\pi}{2}$, is only legitimate if you first insert a regulator such as $\eu^{-sx}$, do the interchange for $s>0$, and let $s\to0^{+}$ with an Abel-summation argument at the end. A naive numerical quadrature is equally unsafe: truncating at $R$ leaves a tail of size $\sim\tfrac{1}{R}$, not $\sim\eu^{-R}$.
\end{trap}

\begin{seealsob}Comparison and limit comparison; the Dirichlet integral, $\Si$ and $\Ci$; Fresnel integrals; justifying the swap.\end{seealsob}

The practical consequence of conditional convergence is a rule about what you are allowed to do next, and it is worth stating as a checklist because the temptation to forget it is strong. With absolute convergence you may expand the integrand in a series and integrate termwise, swap the integral with a limit or a sum, swap the order of a double integral, and differentiate under the integral sign with only a local dominating bound to check. With conditional convergence you may do none of those without a separate argument. The standard repair is to insert a regulator that restores absolute convergence, perform the manipulation, and remove the regulator at the end: replace $\int_0^{\infty}\tfrac{\sin x}{x}\dd x$ by $\int_0^{\infty}\eu^{-sx}\tfrac{\sin x}{x}\dd x=\arctan\tfrac1s$, which is absolutely convergent for $s>0$ and legal to differentiate, then let $s\to0^{+}$ to recover $\tfrac{\pi}{2}$. The regulator is doing exactly the job that the $\varepsilon$ did in the endpoint discipline: it makes the object well defined while the algebra happens.

Conditional convergence is the mildest form of misbehaviour: the integral exists, it merely needs careful handling. The next two sections deal with objects that do not exist at all, and with the two conventions under which they are nevertheless assigned numbers.

\section{Assigning a value to a divergent object}

There are exactly two respectable ways to produce a finite number from an integral that diverges, and both work by imposing a cancellation that the integral itself does not have. The first, the Cauchy principal value, cancels a singularity against its own mirror image by approaching it symmetrically. The second cancels one divergence against another divergence elsewhere in the same expression. Both are conventions, and the obligation in both cases is identical: say which one you used.

Neither convention is arbitrary, and it is worth knowing why, because the reason is what licenses their use in applications. The principal value is the unique symmetric regularisation, and it is the one that arises naturally from contour integration: indenting a contour around a simple pole on the real axis with a small semicircle contributes $\pm\iu\pi\Res$, and the remaining real integral is exactly the principal value. It is also the definition under which the Hilbert transform, the Kramers-Kronig relations, and the Sokhotski-Plemelj formula $\dfrac{1}{x\mp\iu0}=\PV\dfrac1x\pm\iu\pi\delta(x)$ are true. The matched-pair regularisation is likewise not a choice but a bookkeeping device: the object being computed is genuinely convergent, and the divergent halves are an artefact of how it was written. That distinction matters. In the first case you are reporting a new object; in the second you are reporting the original integral, computed carefully.

\tech[Cauchy principal value]{Cauchy principal value}{hard}
\cue A simple pole strictly inside the interval, so that the integrand behaves like $\dfrac{A}{x-c}$ near $c$; or an integral over $\R$ of something like $\dfrac{x}{1+x^{2}}$ that decays like $\tfrac1x$ with opposite signs at the two ends.
\method Approach the singularity symmetrically and take one limit rather than two:
\[\PV\int_a^bf(x)\dd x=\lim_{\varepsilon\to0^{+}}\Bigl[\int_a^{c-\varepsilon}f(x)\dd x+\int_{c+\varepsilon}^{b}f(x)\dd x\Bigr],\qquad \PV\int_{-\infty}^{\infty}f=\lim_{R\to\infty}\int_{-R}^{R}f.\]
The two cut-offs must shrink at the \emph{same rate}, and that single requirement is what makes the odd part of the singularity cancel. A principal value exists for a simple pole (residue-type singularity) and for a $\tfrac1x$ tail; it does not exist for a double pole, because $\tfrac{1}{(x-c)^{2}}$ is even about $c$ and the two sides reinforce instead of cancelling.

\begin{worked}
\[\PV\int_0^{2}\frac{\dd x}{x-1}=\lim_{\varepsilon\to0^{+}}\Bigl[\bigl[\ln\abs{x-1}\bigr]_0^{1-\varepsilon}+\bigl[\ln\abs{x-1}\bigr]_{1+\varepsilon}^{2}\Bigr]=\lim_{\varepsilon\to0^{+}}\bigl[\ln\varepsilon+(0-\ln\varepsilon)\bigr]=0.\]
The ordinary integral does not exist: $\int_0^{1}\tfrac{\dd x}{x-1}=-\infty$ and $\int_1^{2}\tfrac{\dd x}{x-1}=+\infty$. The principal value is $0$, and it is $0$ because the two infinities are the same infinity approached from opposite sides.

The tail version: $\PV\displaystyle\int_{-\infty}^{\infty}\frac{x}{1+x^{2}}\dd x=\lim_{R\to\infty}\bigl[\tfrac12\ln(1+x^{2})\bigr]_{-R}^{R}=0$, while $\int_0^{\infty}\tfrac{x}{1+x^{2}}\dd x=+\infty$ on its own.
\end{worked}

\begin{worked}[A trigonometric principal value]
\[\PV\int_0^{\pi}\tan x\dd x=0.\]
The singularity is at $x=\tfrac{\pi}{2}$, interior, and simple: near it $\tan x\sim\dfrac{-1}{x-\pi/2}$. With symmetric cut-offs,
\[\int_0^{\pi/2-\varepsilon}\tan x\dd x=-\ln\cos\Bigl(\frac{\pi}{2}-\varepsilon\Bigr)=-\ln\sin\varepsilon,\qquad \int_{\pi/2+\varepsilon}^{\pi}\tan x\dd x=\ln\sin\varepsilon,\]
and the two cancel exactly for every $\varepsilon$, so the limit is $0$. Note that they cancel \emph{before} the limit, which is the signature of a clean principal value: no residual constant is left behind. The ordinary integral, of course, does not exist.
\end{worked}

\begin{trap}
Two errors, and the second is subtle. First, reporting a principal value as if it were a convergent integral. It is a different object, and it is not additive over subintervals in the way an integral is; a numerical routine that happens to straddle the pole symmetrically will ``confirm'' it, which proves nothing. Second, the symmetry requirement is not cosmetic. Cutting at $1-\varepsilon$ on the left and $1+2\varepsilon$ on the right in the example above gives $\ln\varepsilon-\ln(2\varepsilon)=-\ln2$, not $0$, and every rate you choose gives a different answer. The principal value is the number attached to the \emph{symmetric} rate, and that choice must be stated.
\end{trap}

\begin{seealsob}Endpoint limit discipline; regularising a matched pair of divergences; disguised zero; quotable residue results.\end{seealsob}

The principal value cancels a singularity against itself. The richer situation, and the one that produces most of the famous constants, is an integrand that is visibly a difference of two divergent pieces, each of which you can integrate exactly. There the cancellation is between two different divergences, and the technique is to make them share a cut-off so the infinite parts can be subtracted before the limit is taken.

\tech[Regularising a matched pair of divergences]{Regularising a matched pair of divergences}{hard}
\cue An integrand that is a difference of two pieces each of which diverges on its own: $\dfrac{1}{\ln x}+\dfrac{1}{1-x}$ on $[0,1]$, $\sec x-\tan x$ at $\tfrac{\pi}{2}$, $\dfrac{f(ax)-f(bx)}{x}$ at both ends, $\dfrac{1-x^{a-1}}{1-x}$ at $1$.
\method Introduce a common cut-off, integrate each divergent half exactly, and watch the $\ln\varepsilon$ terms cancel before letting $\varepsilon\to0$:
\[\int_a^{b}(u-v)=\lim_{\varepsilon\to0^{+}}\Bigl[\int_a^{b-\varepsilon}u-\int_a^{b-\varepsilon}v\Bigr],\]
with the \emph{same} $\varepsilon$ in both. The parametric alternative is often cleaner: introduce a regulator $s$ that makes each half converge, evaluate as an analytic function of $s$, and take the limit of the difference. The finite remainder is where the interesting constant lives.

\begin{worked}
\[\int_0^{1}\Bigl(\frac{1}{\ln x}+\frac{1}{1-x}\Bigr)\dd x=\gamma\approx0.577216.\]
Each half diverges logarithmically at $x=1$: near $1$, $\ln x\sim-(1-x)$, so $\dfrac{1}{\ln x}\sim\dfrac{-1}{1-x}$ and the two terms cancel to leading order, leaving a bounded integrand. Cut both at $1-\varepsilon$. The second half is $-\ln\varepsilon$ exactly. The first, after $x=\eu^{-t}$, is $-\int_{\varepsilon'}^{\infty}\dfrac{\eu^{-t}}{t}\dd t$ with $\varepsilon'\sim\varepsilon$, whose expansion is $\ln\varepsilon+\gamma+O(\varepsilon)$ up to sign. The logarithms cancel and $\gamma$ is the residue. Quadrature on the combined integrand returns $0.5772156649$, which is $\gamma$ to ten places.
\end{worked}

\begin{worked}[The same cancellation, visible in one line]
\[\int_0^{\pi/2}(\sec x-\tan x)\dd x.\]
Both $\int\sec$ and $\int\tan$ diverge at $\tfrac{\pi}{2}$. But $\sec x-\tan x=\dfrac{1-\sin x}{\cos x}$, and the antiderivatives combine: $\ln\abs{\sec x+\tan x}+\ln\abs{\cos x}=\ln\abs{1+\sin x}$. Hence the value is $\bigl[\ln(1+\sin x)\bigr]_0^{\pi/2}=\ln2\approx0.693147$. Here the regularisation is invisible because the two logarithms were combined symbolically before the endpoint was reached, which is the cheapest form of the technique: do the subtraction inside the antiderivative.

The parametric form, for reference: $\displaystyle\int_0^{1}\frac{1-x^{a-1}}{1-x}\dd x=\dig(a)+\gamma$. At $a=3$ the integrand is $1+x$ and the value is $\tfrac32$, matching $\dig(3)+\gamma=\tfrac32$.
\end{worked}

\begin{trap}
Cutting the two halves off at different points. If you truncate $\int\dfrac{\dd x}{\ln x}$ at $1-\varepsilon$ and $\int\dfrac{\dd x}{1-x}$ at $1-2\varepsilon$, the surviving $\ln2$ contaminates the answer, and since $\ln2$ is exactly the size of a plausible constant, the result looks reasonable. This is the mechanism behind almost every ``off by $\ln2$'' error in this family. The same applies to Frullani: $\int_0^{\infty}\dfrac{f(ax)-f(bx)}{x}\dd x=\bigl(f(0)-f(\infty)\bigr)\ln\dfrac ba$ requires the two pieces to be cut at the same $x$, not at the same argument of $f$.
\end{trap}

\begin{seealsob}Frullani's integral; the Euler-Mascheroni constant; Cauchy principal value; conjugate multiplication.\end{seealsob}

\section{Engineered endpoints}

The rest of this chapter is about a different kind of reading. In every technique so far the bounds were given by the problem and the work was in the integrand. There is a large family of problems where the reverse is true: the integrand is routine, and the bounds have been chosen with such deliberate strangeness that they constitute an instruction. A limit of integration written as $\sqrt{\eu^{2}-1}$, or $\log_{8}\eu$, or $\ln^{2}(2\eu)$, is not a number the problem needed. It is a number that becomes clean under one particular substitution, and identifying which one is faster than any inspection of the integrand.

\begin{keynote}[Read the bound as a message]
Ask, of any strange bound $\beta$: which standard substitution sends $\beta$ to something memorable? $\beta=\sqrt{\eu^{2}-1}$ satisfies $1+\beta^{2}=\eu^{2}$, so $u=1+x^{2}$ lands on $\eu^{2}$ and $\ln u$ lands on $2$. $\beta=\log_{8}\eu=\tfrac{1}{\ln8}$ satisfies $\tfrac1\beta=\ln8$, so $u=\tfrac1x$ lands on $\ln8$ and $\eu^{1/x}$ lands on $8$. $\beta=\ln(\pi/4)$ satisfies $\eu^{\beta}=\tfrac{\pi}{4}$, so $u=\eu^{x}$ lands on a special angle. $\beta=\ln^{2}(2\eu)$ satisfies $\sqrt\beta=1+\ln2$, so $u=\sqrt x$ lands on $\ln(2\eu)$ and $\eu^{u}$ lands on $2\eu$. In each case the substitution is then confirmed by the integrand, but it was proposed by the bound.
\end{keynote}

\tech{Endpoint engineering}{easy}
\cue Bounds written strangely when a plain number would have served: $\sqrt{\eu^{2}-1}$, $\log_{8}\eu$ and $\log_{2}\eu$, $\ln^{2}(2\eu)$, $\ln(\pi/4)$ to $\ln(\pi/2)$, $\tan^{-1}3$, $\sqrt{3}/3$. The oddness of the bound is the hint, and it is usually the only hint.
\method Ask what the candidate substitution does to \emph{that number}, and transform the bounds symbolically. Never evaluate a transformed bound numerically before simplifying: $0.4809$ tells you nothing and $\ln8$ tells you everything. If the transformed bounds are clean and the transformed integrand is routine, you have found the intended route; if the bounds go clean and the integrand does not, you have found the wrong substitution and should try the next candidate on the same bound.

\begin{worked}
\[\int_{\log_{8}\eu}^{\log_{2}\eu}\frac{\eu^{1/x}}{x^{2}\bigl(\eu^{1/x}+1\bigr)}\dd x.\]
The bounds are $\tfrac{1}{\ln8}$ and $\tfrac{1}{\ln2}$, screaming for $u=\tfrac1x$. Then $\dd u=-\tfrac{\dd x}{x^{2}}$, the bounds become $\ln8$ and $\ln2$ with the orientation reversed, and
\[I=\int_{\ln2}^{\ln8}\frac{\eu^{u}}{\eu^{u}+1}\dd u=\bigl[\ln(1+\eu^{u})\bigr]_{\ln2}^{\ln8}=\ln9-\ln3=\ln3\approx1.098612.\]
Every piece of the difficulty was in reading $\log_{8}\eu$ as $\tfrac{1}{\ln8}$. Quadrature on the original confirms $1.0986122887$.
\end{worked}

\begin{worked}[Three more bounds, decoded]
\[\int_0^{\sqrt{\eu^{2}-1}}\frac{x\ln(1+x^{2})}{1+x^{2}}\dd x.\]
The bound says $u=1+x^{2}$, which runs from $1$ to $\eu^{2}$. With $\dd u=2x\dd x$,
\[I=\frac12\int_1^{\eu^{2}}\frac{\ln u}{u}\dd u=\frac14\bigl[\ln^{2}u\bigr]_1^{\eu^{2}}=\frac{2^{2}}{4}=1.\]
Next, $\displaystyle\int_{\ln(\pi/4)}^{\ln(\pi/2)}\eu^{x}\cot(\eu^{x})\dd x$: the bounds say $u=\eu^{x}$, running from $\tfrac{\pi}{4}$ to $\tfrac{\pi}{2}$, so $I=\int_{\pi/4}^{\pi/2}\cot u\dd u=\bigl[\ln\sin u\bigr]_{\pi/4}^{\pi/2}=0-\ln\tfrac{1}{\sqrt2}=\tfrac12\ln2\approx0.346574$.

Last, $\displaystyle\int_1^{\ln^{2}(2\eu)}\frac{\eu^{\sqrt x}}{\sqrt x}\dd x$: the squared logarithm says $u=\sqrt x$, running from $1$ to $\ln(2\eu)=1+\ln2$, and $\dd x=2u\dd u$ cancels the $\tfrac1{\sqrt x}$. So $I=2\int_1^{1+\ln2}\eu^{u}\dd u=2\bigl(\eu^{1+\ln2}-\eu\bigr)=2(2\eu-\eu)=2\eu\approx5.436564$.
\end{worked}

\begin{worked}[An inverse-trigonometric bound]
\[\int_0^{\arctan3}\frac{\sec^{2}x}{(1+\tan x)^{2}}\dd x.\]
A bound of $\arctan3$ is an instruction to set $u=\tan x$, which sends it to $3$; the $\sec^{2}x\dd x$ in the numerator is then exactly $\dd u$. So
\[I=\int_0^{3}\frac{\dd u}{(1+u)^{2}}=\Bigl[\frac{-1}{1+u}\Bigr]_0^{3}=1-\frac14=\frac34.\]
The same reading applies to $\arcsin$, $\operatorname{arsinh}$ and $\operatorname{arcosh}$ bounds: each names the value that its own inverse substitution lands on, and the integrand should then contain the matching derivative factor. If it does not, you have misread the bound.
\end{worked}

\begin{trap}
Evaluating the transformed bound as a decimal instead of a symbol. The whole design of these problems is that the transformed bound is a recognisable constant, and a decimal destroys the recognition; $\ln9-\ln3$ is instantly $\ln3$, while $2.1972-1.0986$ is a subtraction you may or may not notice collapses. The second failure is trusting the bound over the integrand: a clean bound under a substitution that leaves the integrand worse is a coincidence, not a signal, and you should back out rather than push on. And when the substitution is not injective on the interval (as $u=x^{2}$ is not on an interval containing $0$), transformed bounds can coincide and silently annihilate the integral.
\end{trap}

\begin{seealsob}Substitution by spotting the differential; exponent and logarithm law collapse; perfect-derivative disguise.\end{seealsob}

The habit generalises past the specific bounds listed above. Whenever a problem contains a number that is more complicated than it needs to be, treat the complication as a specification. A coefficient of $\tfrac{\sqrt3}{3}$ rather than $0.577$ is $\tan\tfrac{\pi}{6}$; a bound of $\tfrac{1}{\sqrt2}$ against an $\arcsin$ is $\sin\tfrac{\pi}{4}$; an upper limit of $\eu^{\pi}$ against a $\ln$ is a request for $\pi$; a limit of $\ln(1+\sqrt2)$ is $\operatorname{arsinh}1$. In every case the author has written the answer to ``where does the substitution land'' into the problem statement, because there is no other reason to write the number that way. This reading takes ten seconds and frequently determines the entire solution, which makes it the highest-yield habit in the chapter.

Endpoint engineering has a converse worth stating, because it closes the loop with the first half of this chapter. Sometimes the strange bound is strange because it is the point where the integrand becomes improper, and the design is not a substitution at all but a test of whether you notice. A bound of $1$ against a $\sqrt{1-x^{2}}$ downstairs, a bound of $0$ against $\ln x$, a bound of $\tfrac{\pi}{2}$ against $\tan x$: these are engineered too, and what they engineer is a limit. The reading protocol is the same in both cases. Look at the bound, ask what the integrand does there, and answer before integrating.

\section{Choosing the route}

The decision procedure for an improper integral is short, and it is short because the first two steps are forced. Locate the bad points; classify each one. Only then does the integrand's structure matter.

\begin{center}
\begin{tikzpicture}[node distance=5mm and 7mm]
\node[dtest] (a) {Where does the integrand\\ or the interval misbehave?};
\node[dtest, below left=9mm and 6mm of a] (b) {Is $\abs f$ integrable\\ near each bad point?\\ ($p$-test, comparison)};
\node[dnode, below right=9mm and 2mm of a] (c) {No bad point:\\ bounds merely strange};
\node[dleaf, below left=9mm and 0mm of b] (d) {Yes: converges absolutely.\\ Evaluate freely; swaps legal};
\node[dtest, below right=9mm and 0mm of b] (e) {Does it oscillate\\ with bounded primitive?};
\node[dleaf, below=9mm of e] (f) {Yes: conditional. Parts once,\\ or alternating series. No DCT.\\ No: divergent, so report\\ $\PV$ or a regularised pair};
\node[dleaf, below=9mm of c] (g) {Decode the bound:\\ which substitution\\ makes it clean?};
\draw[dedge] (a) -- node[dlab,left]{some} (b);
\draw[dedge] (a) -- node[dlab,right]{none} (c);
\draw[dedge] (b) -- node[dlab,left]{yes} (d);
\draw[dedge] (b) -- node[dlab,right]{no} (e);
\draw[dedge] (e) -- (f);
\draw[dedge] (c) -- (g);
\end{tikzpicture}
\end{center}

\begin{keynote}[Making an improper integral safe for a machine]
Numerical quadrature is a poor judge of improperness, so transform before you compute. An endpoint singularity of the form $x^{-p}$ is removed by $x=u^{1/(1-p)}$, or generically by the tanh-sinh map $x=\tanh(\tfrac\pi2\sinh t)$, which pushes both endpoints to $t=\pm\infty$ and makes the transformed integrand decay doubly exponentially. An infinite range is pulled to a finite one by $x=\tfrac{u}{1-u}$ or $x=\eu^{t}$. A conditionally convergent oscillatory tail should be split at the zeros of the oscillating factor and summed as an alternating series, with the tail bounded by the first omitted hump rather than by a truncation guess. Applying these before quadrature is not a refinement; without them the returned digits are unrelated to the integral.
\end{keynote}

Two habits make the tree unnecessary in most cases. The first is to solve for the zeros of every denominator and every logarithm argument on the closed interval before doing anything else, which takes a few seconds and catches the interior pole that no amount of algebraic fluency will catch later. The second is to write down, for each bad point, the single power $x^{-p}$ that matches the local behaviour, and to read the verdict off the $p$-test. Two powers and a conjunction is the entire convergence analysis of most integrals worth writing down.

\begin{keynote}[What to report]
A convergent integral gets a number. A conditionally convergent one gets a number and a sentence saying that it is conditional, because that sentence is what licenses or forbids the reader's next manipulation. A divergent one gets the word ``diverges'', optionally followed by a principal value \emph{explicitly labelled} as such. The one thing that is never acceptable is a bare number for a divergent object, and the reason it is never acceptable is that such a number is usually correct-looking: $-2$ for $\int_{-1}^{1}x^{-2}\dd x$, $0$ for $\int_{-1}^{1}\tfrac{\dd x}{x}$, $\ln2$ for a mismatched regularisation. Wrong answers in this chapter do not look wrong.
\end{keynote}

\section{Recognition matrix}
\begin{recog}{@{}p{0.30\textwidth}p{0.36\textwidth}p{0.28\textwidth}@{}}
\rowcolor{mist}\mh{If you see} & \mh{Reach for} & \mh{If that fails} \\
\midrule
\endhead
An infinite bound & $\lim_{R\to\infty}\int_a^{R}$; test decay against $x^{-p}$, $p>1$ & Substitute $x=1/u$ to move $\infty$ to $0$ \\
$\ln x$ or $x^{-p}$ at an endpoint & $\lim_{\varepsilon\to0^{+}}$; test $p<1$ at the origin & $x=\eu^{-t}$ turns $0$ into $\infty$ \\
A denominator zero inside $[a,b]$ & Split at $c$; both halves must converge & Simple pole: report a $\PV$ \\
$\dfrac{\sin x}{x}$ at $0$, $\dfrac{\ln x}{1-x}$ at $1$ & Removable: the integral is proper & Expand in series to confirm \\
$\dfrac{x^{s-1}}{1+x^{n}}$ on $(0,\infty)$ & Two $p$-tests: $s>0$ and $s<n$ & Value $\dfrac{\pi}{n\sin(\pi s/n)}$ on the strip \\
$\dfrac{1}{x\ln^{q}x}$ near $\infty$ & Log $p$-test: converges iff $q>1$ & $u=\ln x$ \\
$\sqrt{x^{3}+1}$ or similar downstairs & Limit comparison with the leading power & Strip constants, keep the exponent \\
$\dfrac{\sin x}{x}$, $\dfrac{\cos x}{\sqrt x}$ on $(0,\infty)$ & Dirichlet: conditional convergence & Parts once to gain $x^{-2}$ \\
$\sin(x^{2})$, $\cos(x^{2})$ & $u=x^{2}$; decay hides in the Jacobian & Fresnel values $\tfrac12\sqrt{\pi/2}$ \\
$\dfrac{A}{x-c}$ with $c$ interior & Symmetric cut-off: $\PV$ exists & Double pole: no $\PV$ \\
$\dfrac{x}{1+x^{2}}$ on $\R$ & $\PV=0$ by odd symmetry of the tail & State that it is a $\PV$ \\
$\dfrac{1}{\ln x}+\dfrac{1}{1-x}$ & Common cut-off; logs cancel; answer $\gamma$ & Parametric regulator $s$ \\
$\sec x-\tan x$ at $\tfrac{\pi}{2}$ & Combine antiderivatives first: $\ln(1+\sin x)$ & Weierstrass $t=\tan\tfrac x2$ \\
$\dfrac{f(ax)-f(bx)}{x}$ & Frullani with a shared cut-off & Parameter differentiation \\
$\dfrac{1-x^{a-1}}{1-x}$ on $[0,1]$ & $\dig(a)+\gamma$ & Termwise, then telescope \\
Bound $\sqrt{\eu^{2}-1}$ & $u=1+x^{2}$ lands on $\eu^{2}$ & Then $\ln u$ lands on $2$ \\
Bound $\log_{b}\eu$ & $u=1/x$ lands on $\ln b$ & Then $\eu^{1/x}$ lands on $b$ \\
Bound $\ln(\pi/4)$ or $\ln(\pi/2)$ & $u=\eu^{x}$ lands on a special angle & Check the integrand carries $\eu^{x}$ \\
Bound $\ln^{2}(2\eu)$ & $u=\sqrt x$ lands on $1+\ln2$ & Then $\eu^{u}$ lands on $2\eu$ \\
A bound that is exactly a singularity & It is a limit, not a substitution & Classify, then take the limit \\
A clean answer from a blind antiderivative & Check for an interior pole & Sign test: positive integrand, negative value \\
\end{recog}
